\section*{Chapter \ref{ch:b1:water-small-molecules} --- Water and Small Biomolecules}

\begin{solution}{exo:b1:water-small-molecules:1}
The bent, \omterm{def:b1:water-small-molecules:water}{polar molecule} has a partially negative oxygen and partially
positive hydrogens: it surrounds a cation with its oxygens and an anion
with its hydrogens, stabilising them (hydration), and it makes \omterm{def:b1:water-small-molecules:water}{hydrogen
bonds} with polar groups. Oil has no charges or polar groups to bond to;
water molecules around it must order into a cage, which is unfavourable,
so oil is excluded.
\end{solution}

\begin{solution}{exo:b1:water-small-molecules:2}
\omterm{def:b1:water-small-molecules:water}{Hydrogen bond} (\qtyrange{4}{20}{kJ/mol}): \omterm{def:b1:water-small-molecules:ph}{base} pairing in DNA, protein
secondary structure. \omterm{def:b1:water-small-molecules:bonds}{Ionic bond} (about \qty{12}{kJ/mol} in water): salt
bridges in proteins, substrate binding. \omterm{prop:b1:water-small-molecules:properties}{Hydrophobic effect} (not a bond,
but comparable in total): membrane assembly, protein folding. \omterm{def:b1:water-small-molecules:bonds}{Van der
Waals} (\qtyrange{1}{4}{kJ/mol} each): the close packing of any two
complementary surfaces, enzyme--substrate fit.
\end{solution}

\begin{solution}{exo:b1:water-small-molecules:3}
$\mathrm{pH} = -\log(4\times 10^{-8}) = 7.4$. At \omterm{def:b1:water-small-molecules:ph}{pH} 1.5,
$[\mathrm{H^+}] = 10^{-1.5} = \qty{0.032}{mol/L}$.
\end{solution}

\begin{solution}{exo:b1:water-small-molecules:4}
From about 5.8 to 7.8 (p$K_a \pm 1$). At 7.4: $\log(\text{base}/\text{acid})
= 0.6$, ratio $4$.
\end{solution}

\begin{solution}{exo:b1:water-small-molecules:5}
At \omterm{def:b1:water-small-molecules:ph}{pH} 4.76: \qty{50}{\%}. At 7.2: $\log(\mathrm{A^-}/\mathrm{HA}) =
2.44$, ratio $275$: \qty{99.6}{\%} ionised. At \omterm{def:b1:water-small-molecules:ph}{pH} 2: ratio $10^{-2.76} =
1/575$: \qty{0.2}{\%} ionised. The neutral acid form crosses the \omterm{def:b1:membranes-transport:membrane}{lipid
bilayer}; acetic acid is absorbed in the stomach, where it is
un-ionised, and trapped as acetate in the \omterm{def:b1:eukaryotic-cell:organelle}{cytosol}.
\end{solution}

\begin{solution}{exo:b1:water-small-molecules:6}
\omterm{def:b1:water-small-molecules:ph}{pH} $= 6.8 + \log 1 = 6.8$. After HCl: \omterm{def:b1:water-small-molecules:ph}{base} $40$, acid $60$: \omterm{def:b1:water-small-molecules:ph}{pH} $= 6.8 +
\log(40/60) = 6.62$. In pure water, $[\mathrm{H^+}] = \qty{0.01}{mol/L}$:
\omterm{def:b1:water-small-molecules:ph}{pH} $2$.
\end{solution}

\begin{solution}{exo:b1:water-small-molecules:7}
$500/2.4 = \qty{208}{g}$. To vaporise a molecule all four of its \omterm{def:b1:water-small-molecules:water}{hydrogen
bonds} must be broken; \qty{2.4}{kJ/g} is \qty{43}{kJ/mol}, roughly the
energy of two bonds per molecule (each bond being shared by two).
\end{solution}

\begin{solution}{exo:b1:water-small-molecules:8}
$\Delta G^{\circ\prime} = -RT\ln 19 = -2.48\times 2.94 =
\qty{-7.3}{kJ/mol}$. In the \omterm{def:b1:cell-unit-of-life:cell}{cell}: $\Delta G = -7.3 + 2.48\ln(2/0.02) =
-7.3 + 11.4 = \qty{+4.1}{kJ/mol}$: the reaction runs backward, toward
glucose-1-phosphate (as it does in glycogen synthesis).
\end{solution}

\begin{solution}{exo:b1:water-small-molecules:9}
Water is densest at \qty{4}{\celsius}; colder water and ice are lighter
and stay at the surface, so the ice forms on top and insulates the
liquid below, which stays near \qty{4}{\celsius} all winter. If ice
sank, lakes would freeze solid from the bottom up and thaw only at the
surface in summer; freshwater life through a winter would be
impossible.
\end{solution}

\begin{solution}{exo:b1:water-small-molecules:10}
$RT = \qty{2.58}{kJ/mol}$. Rested: $Q = 0.9\times 10^{-3}\times 8\times
10^{-3}/8\times 10^{-3} = 9\times 10^{-4}$, $\ln Q = -7.0$: $\Delta G =
-30.5 - 18.1 = \qty{-48.6}{kJ/mol}$. Fatigued: $Q = 3\times 10^{-3}\times
30\times 10^{-3}/5\times 10^{-3} = 0.018$, $\ln Q = -4.0$: $\Delta G =
-30.5 - 10.4 = \qty{-40.9}{kJ/mol}$. Fatigue has cut the energy per \omterm{def:b1:water-small-molecules:atp}{ATP}
by \qty{16}{\%}, enough to slow the pumps and motors that need most of
it.
\end{solution}

\begin{solution}{exo:b1:water-small-molecules:11}
$[\mathrm{H^+}] = 10^{-7.2} = \qty{6.3e-8}{mol/L}$; volume
\qty{1e-12}{L}: $\qty{6.3e-20}{mol}\times\num{6e23} = 38$ protons. Against
$\num{3e9}$ buffering groups: the ``\omterm{def:b1:water-small-molecules:ph}{pH}'' is not a count of protons but the
state of protonation of the \omterm{def:b1:water-small-molecules:buffer}{buffers}, which exchange protons with the
water a billion times faster than the few free ones could matter.
\end{solution}

\begin{solution}{exo:b1:water-small-molecules:12}
At equilibrium \omterm{def:b1:water-small-molecules:atp}{ATP}, ADP and phosphate would sit at a ratio of about
$10^{-5}$; the \omterm{def:b1:cell-unit-of-life:cell}{cell} holds \omterm{def:b1:water-small-molecules:atp}{ATP} a thousand times above ADP, so that its
\omterm{prop:b1:water-small-molecules:families}{hydrolysis} releases \qty{50}{kJ/mol} rather than nothing. Every uphill
process --- synthesis, transport, movement --- is coupled to that
displacement; catabolism's only role is to maintain it. A \omterm{def:b1:cell-unit-of-life:cell}{cell} whose \omterm{def:b1:water-small-molecules:atp}{ATP}
reaches equilibrium has no gradient to spend: pumps stop, ions leak,
the \omterm{def:b1:cell-unit-of-life:cell}{cell} swells and dies within minutes. Being alive is the maintenance
of that disequilibrium, not a substance.
\end{solution}

\begin{solution}{pb:b1:water-small-molecules:1}
\textbf{1.} $6.1 + \log(24/1.2) = 6.1 + 1.301 = 7.40$.
\textbf{2.} $10^{-7.4} = \qty{4.0e-8}{mol/L} = \qty{40}{nmol/L}$.
\textbf{3.} $40\times 5 = \qty{200}{nmol}$, six hundred thousand times
fewer than the bicarbonate.
\textbf{4.} $\mathrm{CO_2} = 1.2\times 46/40 = \qty{1.38}{mmol/L}$;
\omterm{def:b1:water-small-molecules:ph}{pH} $= 6.1 + \log(25/1.38) = 6.1 + 1.258 = 7.36$.
\textbf{5.} In a flask the ratio $20:1$ \omterm{def:b1:flowering-plant-organization:organs}{leaves} little acid partner to
absorb \omterm{def:b1:water-small-molecules:ph}{base}, and the capacity is low at one unit from the p$K_a$; in
the body the acid partner is a gas whose concentration the lungs fix
and can regenerate without limit, so the ratio can be reset by
ventilation regardless of the p$K_a$.
\textbf{6.} $24/1.2 = 20$; the \omterm{def:b1:water-small-molecules:ph}{base} is $20/21 = \qty{95}{\%}$ of the pair.
\textbf{7.} $60/5 = \qty{12}{mmol/L}$.
\textbf{8.} Bicarbonate $24 - 12 = \qty{12}{mmol/L}$; $\mathrm{CO_2}$
$1.2 + 12 = \qty{13.2}{mmol/L}$.
\textbf{9.} \omterm{def:b1:water-small-molecules:ph}{pH} $= 6.1 + \log(12/13.2) = 6.1 - 0.04 = 6.06$.
\textbf{10.} $[\mathrm{H^+}] = \qty{0.012}{mol/L}$: \omterm{def:b1:water-small-molecules:ph}{pH} $1.92$.
\textbf{11.} A shift of $1.34$ units: it saves the blood from \omterm{def:b1:water-small-molecules:ph}{pH} 1.9
but 6.06 is far below what any enzyme tolerates. Poor, as a closed
system.
\textbf{12.} $[\mathrm{H^+}]$ goes from \qty{40}{nmol/L} to $10^{-6.06} = \qty{870}{nmol/L}$: a rise of \qty{0.83}{\micro mol/L} for \qty{12}{mmol/L} added, i.e.\ one proton in \num{14000} stays free; the rest are on bicarbonate.
\textbf{13.} \omterm{def:b1:water-small-molecules:ph}{pH} $= 6.1 + \log(12/1.2) = 6.1 + 1.0 = 7.10$.
\textbf{14.} The $\mathrm{CO_2}$ generated: $\qty{12}{mmol/L}\times
\qty{5}{L} = \qty{60}{mmol}$.
\textbf{15.} \qty{60}{mmol} in a minute on top of \qty{10}{mmol}:
sevenfold.
\textbf{16.} $\mathrm{CO_2} = \qty{0.9}{mmol/L}$; \omterm{def:b1:water-small-molecules:ph}{pH} $= 6.1 +
\log(12/0.9) = 6.1 + 1.125 = 7.22$.
\textbf{17.} Open bicarbonate near 7.4: the \omterm{def:b1:water-small-molecules:ph}{pH} changes by one unit
when the bicarbonate falls tenfold, from \qty{24}{} to
\qty{2.4}{mmol/L}, i.e.\ \qty{21.6}{mmol/L} of acid per unit; but over
the first units the local slope is $2.3\times[\mathrm{HCO_3^-}] =
\qty{55}{mmol/L}$ per unit. Against the proteins' \qty{25}{mmol/L} per
unit, bicarbonate takes about $55/80 = \qty{69}{\%}$ of the protons,
proteins \qty{31}{\%}.
\textbf{18.} Bicarbonate takes $0.69\times 12 = \qty{8.3}{mmol/L}$:
$24 - 8.3 = \qty{15.7}{mmol/L}$; \omterm{def:b1:water-small-molecules:ph}{pH} $= 6.1 + \log(15.7/1.2) = 6.1 +
1.117 = 7.22$ (the proteins' share consistent with \qty{25}{mmol/L}
per unit $\times 0.18$ units $= \qty{4.5}{mmol/L}$, close to their
\qty{3.7}{mmol/L}).
\textbf{19.} Removing the acid regenerates the bicarbonate consumed
(\qty{12}{mmol/L} back to \qty{24}{}) and, with $\mathrm{CO_2}$ held,
the \omterm{def:b1:water-small-molecules:ph}{pH} returns to 7.40.
\textbf{20.} $60/5 = \qty{12}{mmol/L}$ of bicarbonate lost per day;
after three days $24 - 36 < 0$: bicarbonate is exhausted and the \omterm{def:b1:water-small-molecules:ph}{pH}
collapses well below 7 --- in practice the patient must be given
bicarbonate or dialysed; after one day alone, $6.1 + \log(12/1.2) =
7.10$.
\textbf{21.} $\mathrm{CO_2} = 1.2\times 60/40 = \qty{1.8}{mmol/L}$;
\omterm{def:b1:water-small-molecules:ph}{pH} $= 6.1 + \log(24/1.8) = 6.1 + 1.125 = 7.22$. With bicarbonate at
\qty{33}{}: $6.1 + \log(33/1.8) = 6.1 + 1.263 = 7.36$.
\textbf{22.} The lungs change the denominator (dissolved
$\mathrm{CO_2}$) within breaths, because it is a gas they exhale; the
kidneys change the numerator (bicarbonate) by excreting acid and
regenerating bicarbonate, a process of hours to days.
\textbf{23.} At \omterm{def:b1:water-small-molecules:ph}{pH} 5, $[\mathrm{H^+}] = \qty{10}{\micro mol/L}$: \qty{15}{\micro mol} free in \qty{1.5}{L}, a four-thousandth of the \qty{60}{mmol}; the acid \omterm{def:b1:flowering-plant-organization:organs}{leaves} bound to \omterm{def:b1:water-small-molecules:buffer}{buffers} --- as ammonium and as dihydrogen phosphate.
\textbf{24.} Losing acid raises \omterm{def:b1:water-small-molecules:ph}{pH}: bicarbonate rises by
$100/5 = \qty{20}{mmol/L}$ to \qty{44}{}; \omterm{def:b1:water-small-molecules:ph}{pH} $= 6.1 + \log(44/1.2) =
6.1 + 1.564 = 7.66$ (alkalosis), before the lungs and kidneys
compensate.
\textbf{25.} Closed: from 7.40 to 6.06, a shift of $-1.34$. Open,
bicarbonate alone: to 7.10, a shift of $-0.30$. Open with proteins:
to 7.22, a shift of $-0.18$.
\end{solution}
